\clearpage
\setcounter{page}{1}
\maketitlesupplementary

\section{Datasets}
\label{sec:suppl_datasets}

\subsection{Detection prompt training data}
\label{sec:suppl_detection_data}

For each concept $c$, the positive prompt set $\mathcal B_c^{\mathrm{pos}}$ describes scenes containing 
$c$, and the negative set $\mathcal B_c^{\mathrm{neg}}$ covers both related and unrelated content 
without $c$. 
% The accompanying code provides nudity prompt texts in \texttt{i2p\_detect.py} and detection training templates for IP characters, objects, and artistic styles in \texttt{datasets.py}.

\noindent\textbf{Nudity.}
We train the nudity detection prompt using 960 positive prompts and 480 negative prompts. 
The positive prompts describe full or partial nudity using explicit wording or contextual cues. 
The negatives include clothed people, scenes where clothing or framing conceals the relevant body areas, 
related settings without people, and unrelated scenes.

\noindent\textbf{IP characters and objects.}
We provide 50 shared scene templates. Each specifies a setting and leaves one slot for the subject, as in \emph{At the edge of a courtyard, place \{concept\} where a band of sunlight crosses the paving.} Filling the slot with the target concept gives a positive prompt. Filling it with a related or unrelated concept gives a negative prompt for the same setting. For the Pikachu detector, the example template can describe Pikachu as a positive, Raichu as a related negative, or a teapot as an unrelated negative. These scene templates differ in wording and structure from the IP and object evaluation templates below.

\noindent\textbf{Artistic styles.}
The 17 detection training templates vary how they describe an artwork, as in \texttt{a painting in the style of \{artist\}} and \texttt{a close-up painting in the style of \{artist\}}. We fill the slot with the target artist for positives and with a related artist or an unrelated style for negatives. For Van Gogh, Monet is a related negative and pixel art is an unrelated negative.

\noindent\textbf{TI-flow images.}
For each target concept $c$, we condition the frozen base model on text describing $c$ and generate 100 images to form the fixed image set $\mathcal I_c$. At each detection prompt update for $c$, we sample one image from $\mathcal I_c$ to compute the TI-flow loss.

\subsection{Suppression prompt training data}
\label{sec:suppl_suppression_data}

Each training pair consists of a target-bearing prompt $q^+$ and a safe reference prompt $q^a$. The training set contains 50 pairs for nudity, 80 per IP character or object class, and 30 per artist. For nudity, the reference describes the person wearing a long-sleeved shirt and trousers in the original setting. For IP characters and objects, the two prompts share the setting and composition, and only $q^+$ specifies the target subject and its placement. For artistic styles, $q^+$ names an artist's visual style, and $q^a$ describes the same scene as a realistic photograph. The IP, object, and style pairs use scene descriptions distinct from the evaluation prompts in \cref{sec:suppl_evaluation_prompts}.

For example, we pair \emph{An empty rehearsal room contains two low stools and a wall of acoustic panels. Render the scene in the visual style of Van Gogh.} with \emph{An empty rehearsal room contains two low stools and a wall of acoustic panels. Depict the scene as a realistic photograph.}

\subsection{Evaluation prompts}
\label{sec:suppl_evaluation_prompts}

\noindent\textbf{IP character erasure.}
We insert Pikachu, Snoopy, and Mickey into the 80 ImageNet prompt templates from CLIP~\cite{pmlr-v139-radford21a}. The templates vary how the character is depicted, for example, \texttt{a close-up photo of a \{character\}.} and \texttt{a sculpture of a \{character\}.}

\noindent\textbf{Artistic styles.}
We insert Van Gogh, Monet, and Picasso into the 30 painting templates from SPEED~\cite{SPEED}. Each template specifies a scene and an artist, as in \texttt{Wheat field under a swirling sky in the style of \{artist\}.}

\noindent\textbf{Objects.}
We use \texttt{a photo of \{class\}} for each of the ten ImageNet classes in the object erasure evaluation, yielding prompts such as \texttt{a photo of church} and \texttt{a photo of golf ball}.
